\chapter{Faire évoluer cette documentation}

\section{Principe}

Cette documentation est conçue pour évoluer par petits ajouts. Il vaut mieux
compléter régulièrement un chapitre existant que réécrire tout le manuel à
chaque nouvelle version.

\section{Ajouter une capture}

Placer l'image dans \fichier{docs/latex/figures/}, puis l'insérer avec :

\begin{verbatim}
\begin{figure}[H]
  \centering
  \includegraphics[width=0.9\textwidth]{figures/nom_capture.png}
  \caption{Description de la capture}
\end{figure}
\end{verbatim}

\section{Ajouter un chapitre}

\begin{enumerate}
  \item Créer un fichier dans \fichier{chapters/}.
  \item Ajouter une ligne \commande{\textbackslash input\{chapters/nom\}} dans
  \fichier{main.tex}.
  \item Compiler le document.
\end{enumerate}

\section{Marquer les parties à compléter}

Utiliser la commande :

\begin{verbatim}
\todoDoc{Décrire ici ce qu'il faudra ajouter plus tard.}
\end{verbatim}

\section{Compiler}

Depuis le dossier \fichier{docs/latex} :

\begin{verbatim}
latexmk -pdf main.tex
\end{verbatim}

Si \commande{latexmk} n'est pas disponible :

\begin{verbatim}
pdflatex main.tex
bibtex main
pdflatex main.tex
pdflatex main.tex
\end{verbatim}
